\section{Tidewater grant application}

Tidewater grant application

Lead investigator: Gérard Côté\\
Tidewater Lab, 12 Quay Road\\
Port Mersey, Nova Scotia

\section{What are you working on?}

\emph{Describe the problem in two or three sentences. Plain language is fine.}

Tidewater is a tool that maps river sediment. It turns core samples, depth readings and historical surveys into a map of where silt, sand and clay have moved over the last fifty years. The project began as a pilot in 2025 — small as it was, the pilot showed that the method works on a real estuary.

Our \textbf{main claim} is simple: sediment maps can be cheap. Most \emph{existing} surveys cost \textbf{\emph{far more}} than a town can afford, and the data often sits unused in \texttt{csv} files or in filing cabinets. We use \texttt{tw\_map()} to turn raw depth logs into a layered map that a council engineer can read.

Read more about \href{https://example.org/pilot}{the pilot}. The full data policy is at \url{https://example.org}.

The budget request is \$150,000 over two years. Field costs run about \textasciitilde{}5m per site, and roughly \textasciitilde{}40\% of the budget goes to laboratory work. Software licences are offered at 50 \% off for public bodies.

Column names use snake\_case throughout. Templates use \{braces\} for placeholders, and on Windows the data lives under C:\textbackslash{}path\textbackslash{}tidewater. The model is quadratic in depth, so the fit includes an x\textasciicircum{}2 term. See @ref for the sampling protocol, and wrap values in <tag> when exporting HTML.

Our R\&D partners will review the design. Their comments are welcome — the more critical, the better — and we have built them into the schedule.

The plan is in our own words, not “borrowed” from a funder’s template, and the team’s plan has been checked by two outside readers.

\begin{enumerate}
  \item Collect cores from three sites.
  \item Label each sample by depth and date.
  \item Run the mapping pipeline.
  \begin{itemize}
    \item Check the depth readings against the survey.
    \item Flag any core older than ten years.
  \end{itemize}
  \item Publish the maps under an open licence.
  \item Train two council staff to run the tool.
\end{enumerate}

Items 1-3 are the technical core, and items 4 and 5 are outreach. If \#1 slips, the rest of the schedule slips with it.

\subsection{What would you do with funding?}

\emph{Be specific about what the money buys.}

Funding pays for three kinds of work. Each one maps onto a line of the budget.

\begin{itemize}
  \item Field work
  \begin{enumerate}
    \item Collect cores at each site
    \begin{enumerate}
      \item Label samples on site
    \end{enumerate}
    \item Survey the channel by boat
  \end{enumerate}
  \item Lab work
  \item Software
\end{itemize}

\begin{tabular}{l l l}
Site & Depth (m) & Sample \\
\hline
Lower Tidewater & 4.5 & Silt \\
Mill Creek & 2.0 & Sand \\
\end{tabular}

\begin{quote}
Sediment records where the river has been. Tidewater shows where it is going.
\end{quote}

A sample analysis script looks like this:

\begin{verbatim}
library(tidewater)
cores <- read_cores("data/cores.csv")
map <- tw_map(cores, layer = "silt")
\end{verbatim}

Thank you for reading. We look forward to hearing from you.
